%% kapa: exam=1 round=35 year=2024 period=1 form=A subject=economics subject_ko=경제학원론 seq=18 number=58
%% source: Q-Net 감정평가사 기출문제 (cst00309) artlSeq=5213229 file=raw/1008/5213229/1교시 문제지 원본.pdf
%% conversion: kordoc 4.24.51 -> tools/gen1.py (auto); review=proofread-sonnet; flags=none
\begin{question}{58}
사회후생 관점에서 자원의 효율적 활용에 관한 설명으로 옳지 않은 것은?
\begin{choices}
  \item 계약곡선 상의 점들은 생산의 효율성을 보장하는 점들의 집합이다.
  \item 효용가능곡선은 주어진 상품을 두 사람에게 배분할 때, 두 사람이 얻을 수 있는 최대한의 효용수준의 조합이다.
  \item 효용가능경계란 한 경제에 존재하는 경제적 자원을 가장 효율적으로 배분했을 때 얻을 수 있는 효용수준의 조합이다.
  \item 종합적 효율성(overall efficiency)이란 생산의 효율성과 교환의 효율성이 동시에 달성된 상태를 말한다.
  \item 생산가능곡선은 한 나라의 경제가 주어진 생산요소와 생산기술을 사용하여 최대한 생산할 수 있는 산출물들의 조합이다.
\end{choices}
\end{question}
\answer{1}
